

\chapter{Metodologia pratica} 
\lipsum[1-3]
\section{Aula 1}
O meu trabalho é sobre documentar em \textbf{LaTex} tudo que eu vou aprender no curso de \textbf{LaTex} que estou cursando, nesta primeira aula o professor nos mostrou como foi o processo de criação do \textbf{LaTex}, e quando foi criado, também nos mostrou alguns códigos básicos de \textbf{LaTex} e a ferramenta do \textbf{overleaf} para a criação de documentos em \textbf{LaTex}.


O \textbf{Tex} foi criado em \textbf{1970} por \textbf{Donald Knuth} pois ele queria uma forma mais padronizada de criar os seus livros, em \textbf{1985} \textbf{Leslie Lamport} viu a ideia do \textbf{Tex} e achou uma ideia brilhante e com base no \textbf{Tex} criou uma forma mais fácil e rápida de usar o \textbf{Tex} nomeando a sua forma de \textbf{LaTex} que é usada mundialmente ate hoje.


Foi se destacado na apresentação o uso de paradigmas como \textbf{WYSIWYG} e \textbf{WYSIWYM} que podemos usar o Word como exemplo e o LaTex, o paradigma \textbf{WYSIWYG} significa oque você vê é oque você pega, que no contexto do Word diz que oque você ve é onde você clica, e o \textbf{WYSIWYM} significa oque você vê é oque você quer que diz que você que manda no texto e não onde você clica.


Para colocar \% você precisa escrever \textbackslash \%.

\section{Aula 2}
Na 2º aula do Curso de LaTex nos aprendemos comandos novos de modelo, como Book (Modelo de Livro), Article (Modelo de Artigo), Report (Para artigos mais longos), Também aprendemos comandos de organização de texto, como pular paginas,como dividir em capitulo, seção, subseção e subsubseção, também aprendemos a colocar afiliação aos autores,montamos o chassi do nosso artigo final do curso, e também formas de modificar o modelo (Opções de Classe) como em colunas e outras coisas.


Aprendemos vários modelos de LaTex, desde um artigo curto, ate como fazer slides usando LaTex, e os modelo tem as opções de classe, algumas delas são o twocolumns que faz o texto dividir em 2 colunas em uma pagina apenas, a4paper que deixa os documentos no modelo padrão de documentos, temos o draft que e basicamente uma forma de achar erros de formatação enquanto você ainda esta escrevendo o texto, ele também não carrega as imagens para executar o programa mais rápido e acelera o processamento da compilação.


Na continuação da aula o professor nos ensinou o pacote para colocar margens no texto, que é o geometry, nele você precisa dizer todos os 4 lados como cima, baixo, esquerda e direita, também aprendemos que existem alguns pacotes básicos para o funcionamentos do programa ocorrer de forma correta, como o [T1]{fontenc}, que mapeia a fonte no PDF para quando alguém for copiar o texto conseguir copiar os acentos também, [utf8]{inputenc} que permite que o computador compreenda as formas de acentuação, e o [brazilian]{babel} que faz o computador compreender as  formas de hifenização brasileiras. 


E por final aprendemos o conceitos do abstract, lipsum e o thanks, o abstract nos permite criar uma parte no texto que funciona como um resumo deixando você escrever um texto rápido sobre o artigo ou livro ou oque você for fazer, o lipsum funciona como um criador de texto fantasma, ele cria quantos parágrafos você quiser para ter mais volume de texto, e o thanks serve para filiação dos autores, o comando thanks cria uma parte que mostra informações do autor no final.

\section{Aula 3}

Na aula 3 aprendemos vários tipos diferentes de fontes e como utilizar cada fonte, aprendemos como colocar espaçamento entre linhas e quais espaçamentos, numerar itens como listas, como criar uma referencia a alguma parte no texto no final da pagina, as fontes de LaTeX tem vários tamanhos desde muito pequenos a gigantescos, formas de centralizar o texto na pagina, e como remover a  serifa das bordas das letras ou coloca palavras ou letras sublinhadas.


As fontes no Latex sao muitos variáveis e nelas você modifica o tamanho da fontes usando comandos como "normalsize" ou "Huge", a normalsize serve para tamanhos comuns de fonte, a Huge serve para uma escrita grotescamente grande, e não existem apenas estas duas fontes, existem muito mais para serem utilizadas, existe o espaçamentos entre as linhas que também existe varias formas, porem eles variam entre espaçamento de 2 mm, 1,5 ou 1 e outros O LaTex usa normalmente um espaçamento de linhas de 1.0, porem você pode alterar o espaçamentos importando um pacote direto no preambulo, esse pacote se chama "usepackage {setspace}".


Para finalizar a aula, o professor explicou o conceito de enumeração dos tópicos, oque e muitos útil pra quando você for citar muitos tópicos no seu artigo ou jornal, ele funciona como uma lista que a cada linha adiciona 1 numero no começo, isso e muito útil pra quando as pessoas querem citar o tópico ou ir mais rápido ate ele, e explicou sobre centralização de texto que serve para deixar o texto visualmente mais bonitos para leitura.

\section{Aula 4}
Não consegui comparecer a aula 4. 
\section{Aula 5}

Nesta aula aprendemos a inserir imagens no LaTex e como funciona a inserção de imagens, no LaTex a inserção de fotos no texto e bem parecido com a forma de inserção de fotos do python, uma das partes boas do LaTex comparado ao editor de textos word ou docs é que utilizando IA nesses softwares de edição mais comuns, quando você adiciona uma foto, todo o seu texto fica fora de ordem e tem que passar muitos tempo corrigindo os erros de formatação do texto, no LaTex não, você só precisa escrever algumas linhas de códigos e as imagens ficam corretas no lugar sem atrapalhar a formatação do texto.
\vspace{0.5cm}
\begin{figure}[htbp]
    \centering
    \includegraphics[width=0.1\linewidth]{IMG/Linux_mascot_tux.png}
    \caption{Imagem para exemplo}
    \label{fig:bancada}
\end{figure}

Apos inserir as imagens, temos varias formar de posicionar elas na pagina, como no meio, a esquerda, direita, em cima, em baixo entre outros, e para fazer isso apenas precisa digitar algumas linhas de código, oque facilita muito para o escritor quando estiver criando o artigo,
as mais usadas são: centering (para centralizar o texto na pagina), flushleft(para deixar a imagem a esquerda da pagina), flushright(para deixar a imagem a direita da pagina) entre outras.


Também aprendemos sobre o uso de tabelas no código de forma organizada, as tabelas são muito útil para conter dados de testes ou informações sobre o artigo que esta sendo criado, vou citar uma tabela criada em LaTex e uma imagem abaixo: 

\begin{figure}[htbp]
    \centering
    \includegraphics[width=0.5\linewidth]{IMG/images.png}
    \caption{Imagem para exemplo}
    \label{fig:bancada}
\end{figure}

\begin{table}[htbp]
    \centering
    \caption{Exemplo de tabela}
    \label{tab:exemplo}
    \begin{tabular}{ccc}
        \toprule
        \textbf{Aula1} & \textbf{Aula2} & \textbf{Aula3} \\
        \midrule
        13:30 & 13:30 B1 & 13:30 \\
        17:30 & 17:30 & 17:30 \\
        \bottomrule
    \end{tabular}
\end{table}

\section{Aula 6}
Nesta aula aprendemos sobre citações e referencias cruzadas no LaTex e qual a utilidade delas no texto.
\\Oque são referencias cruzadas?


Referencias cruzadas funcionam como uma marcação, elas servem para caso a forma do texto mude, não trocar a figura, tabela ou citação, em vez de enumerar o documento manualmente você marca o elemento com um identificador e o LaTeX atualiza a numeração automaticamente a cada compilação.
\\Para você usar elas no código, existem 3 formas
\begin{itemize}
    \item \textbackslash label\{marcador\}: Atribui um nome de identificação único ao elemento (como uma seção, figura, tabela ou equação).
    \item \textbackslash ref\{marcador\}: Insere automaticamente o número correspondente ao elemento no texto.
    \item \textbackslash pageref\{marcador\}: Insere o número da página onde o elemento se encontra
\end{itemize}
E agora oque são citações?\\


As citações são formas de repetir algo que outra pessoa já disse, escreveu ou aconteceu em algum momento, filme, livro, artigo etc..., geralmente um artigo tem varias citações de outros artigos que ele se baseou, por isso aprender a fazer citações em LaTex e algo muito importante e interessante, para utilizar uma citação você vai precisar de um arquivo .bib com as informações do conteúdo no qual você esta citando como nome, autor, data, Edição etc..., apos colocar tudo isso no arquivo .bib, você apenas precisa chamar no texto a função \textbackslash cite\{nome que você indicou a citação\}, vou deixar anexado apos este texto um exemplo de uma citação do artigo de lucas Alexandre Zick sobre teleoperação de robôs.   
\cite{zickgames2024teleoperation}
\section{Aula 7}
Na aula 7 aprendemos como utilizar modelos academicos prontos e colocar os dados do nosso artigo dentro deste modelo, modelos são muito uteis e necessarios dentro de um artigo para deixar os artigos de uma revista em um padrão legivel para o leitor, varias revistar tem seus modelos prontos como a Nature, IEEE, Science etc\dots


textbf{abnTeX2}: o abnTeX2 é uma suíte desenvolvida por pesquisadores brasileiros que torna o LaTex uma mao na roda quando se fala em criar trabalhos em ABNT, pois ele fornece estruturas prontas em ABNT, quando ativamos a classe abntex2, ganhamos o acesso a comandos magicos como 
\begin{itemize}
    \item \textbackslash imprimirfolhaderosto
    \item \textbackslash begin\{Dedicatoria\}
    \item \textbackslash begin\{Agradecimentos\}
    \item \textbackslash begin\{Epigrafe\}
\end{itemize}


No Word uma das partes mais demoradas é gerar o indice de figuras ou tabelas. No abnTeX2 quando você utiliza o \textbackslash caption em todas as imagens, o trabalho é nulo, basta apenas usar estes comandos: 
\begin{itemize}
    \item \textbackslash tableofcontents
    \item \textbackslash listoffigures
    \item \textbackslash listoftables
\end{itemize}


A melhor partes da comunidade cientifica é que as revistas fornecem os modelos prontos para você utilizar no seu artigo, o paocte do modelo geralmente vem em um arquivo .zip contendo: um arquivo .cls, .sty, .bst e um main.tex.
\section{Aula 8}
Na aula 8 aprendemos sobre como rodar o LaTex localmente, em vez de pelo overleaf, e quais as vantagens sobre rodar localmente, algumas vantagens de rodar localmente são: o "timeout", ele e um erro que acontece quando você tentar rodar um pdf com muito espaço para armazenar e o compilador online não aceita compilar o seu código, a privacidade e o sigilo tambem sao considerados vantagens quando se compila localmente, pois os dados sensiveis e uma pesquisa nao devem rodar em servidores de terceiros, e a possibilidade de rodar scripts locais que geram tabelas dinamicas direto no LaTex. 


Quando você vai rodar localmente você tem 3 opçoes para rodar o LaTex:

 \begin{itemize}
    \item  \textbf{pdfLaTeX}: essa é a opção mais utilizada, pois é estavel e é o padrão da maioria da revistas, porem nao lida muito bem com fontes modernas, ocasionando a necessidade de pacotes.
    \item \textbf{XeLaTex}: é a opção mais moderna, ela permite que voce digite '\textbackslash setmainfontArial' e o documento completo comece a usar as fontes nativas do seu windows ou linux.
    \item \textbf{LuaLaTex}: E considerado o futuro do LaTex, integrou a linguagem de programação Lua dentro do motor de tipografia, permitindo calculos complexos durante a compilação.
 \end{itemize}   

Rodando localmente o LaTex pode aparecer alguns lixos temporarios, esses lixos Você pode apagar e o codigo vai continuar rodando normalmente, as unicas pastas que sao insubstituíveis são o .tex e a .bib e algumas pastas que podem conter imagens.

\section{Aula 9}


\section{Aula 10}
\lipsum[1-3]
\section{Aula 11}
\lipsum[1-3]
\section{Aula 12}
\lipsum[1-3]


